\element{1.2.1}{
  \begin{questionmult}{ALI}\bareme{haut=2,p=0}
    Pour un ALI idéal, on a
    \begin{multicols}{3}
      \begin{reponses}
        \mauvaise{$\epsilon\approx 0$}
        \bonne{$i_+\approx i_- \approx 0$}
        \mauvaise{$s=\pm V_{\text{sat}}$}
        \mauvaise{$S(p)=\frac{H_0}{1-\tau p}\mathcal{E}(p)$}
        \mauvaise{Le système est nécessairement stable}
      \end{reponses}
    \end{multicols}
  \end{questionmult}
}

\element{1.2.2}{
  \begin{questionmult}{ALIStable}\bareme{haut=2,p=0}
    Pour un ALI idéal, en fonctionnement stable, on a
    \begin{multicols}{3}
      \begin{reponses}
        \bonne{$\epsilon\approx 0$}
        \bonne{$i_+\approx i_- \approx 0$}
        \mauvaise{$s=\pm V_{\text{sat}}$}
        \mauvaise{$S(p)=\frac{H_0}{1-\tau p}\mathcal{E}(p)$}
      \end{reponses}
    \end{multicols}
  \end{questionmult}
}
\element{1.2.2}{
  \begin{questionmult}{ALISnstable}\bareme{haut=2,p=0}
    Pour un ALI idéal, en fonctionnement instable, on a
    \begin{multicols}{3}
      \begin{reponses}
        \mauvaise{$\epsilon\approx 0$}
        \bonne{$i_+\approx i_- \approx 0$}
        \bonne{$s=\pm V_{\text{sat}}$}
        \mauvaise{$S(p)=\frac{H_0}{1-\tau p}\mathcal{E}(p)$}
      \end{reponses}
    \end{multicols}
  \end{questionmult}
}
\setgroupmode{1.2.2}{cyclic}

\element{1.2.3}{
  \begin{questionmult}{stabiliteMontage}\bareme{haut=2,p=0}
    Parmi les montages suivants, lesquels sont stables ?
    \begin{multicols}{2}
      \begin{reponses}
        \bonne{%
          \begin{circuitikz}[scale=0.7,transform shape]%suiveur
            \draw (0,0) node[en amp] (ALI) {};
            \draw (ALI.+) --++ (-1,0) node[left] {$V_e$};
            \draw (ALI.-) -|++ (0,1) -| (ALI.out) --++ (0.5,0) node[right] {$V_s$};
          \end{circuitikz}
        }
        \bonne{%
          \begin{circuitikz}[scale=0.7,transform shape]%ANI
            \draw (0,0) node[en amp] (ALI) {};
            \draw (ALI.+) -|++ (-0,-0.7) node[ground] {};
            \draw (ALI.-) to[R=$R_1$] ++ (-3,0) node[left] {$V_e$};
            \draw (ALI.-) --++ (0,1.5) to[R=$R_2$] ++ (2.5,0) |- (ALI.out) --++ (0.5,0) node[right] {$V_s$};
          \end{circuitikz}
        }
        \bonne{%
          \begin{circuitikz}[scale=0.7,transform shape]%intégrateur
            \draw (0,0) node[en amp] (ALI) {};
            \draw (ALI.+) -|++ (-0,-0.7) node[ground] {};
            \draw (ALI.-) to[R=$R$] ++ (-3,0) node[left] {$V_e$};
            \draw (ALI.-) --++ (0,1.5) to[C=$C$] ++ (2.5,0) |- (ALI.out) --++ (0.5,0) node[right] {$V_s$};
          \end{circuitikz}
        }
        \mauvaise{%
          \begin{circuitikz}[scale=0.7,transform shape]%comparateur
            \draw (0,0) node[en amp,yscale=-1] (ALI) {};
            \draw (ALI.-) --++ (-0.2,0) to[V_<=$V_0$] ++ (0,-1.5) node[ground] {};
            \draw (ALI.+) --++ (-0.7,0) node[left] {$V_e$};
            \draw (ALI.out) --++ (0.7,0) node[right] {$V_s$};
          \end{circuitikz}
        }
        \mauvaise{%
          \begin{circuitikz}[scale=0.7,transform shape]%CH
            \draw (0,0) node[en amp] (ALI) {};
            \draw (ALI.-) node[left] {$V_e$} ;
            \draw (ALI.+) to[R=$R_1$] ++ (-3,0) -|++ (-0,-0.7) node[ground] {};
            \draw (ALI.+) --++ (0,-1.5) to[R=$R_2$] ++ (2.5,0) |- (ALI.out) --++ (0.5,0) node[right] {$V_s$};
          \end{circuitikz}
        }
      \end{reponses}
    \end{multicols}
  \end{questionmult}
}
